% ====================================================================================
\chapter{视觉--语言--动作模型、真机强化学习与展望}
\chaptermark{VLA、真机强化学习与展望}
% ====================================================================================

\section*{"把灭绝的动物捡起来"}

2023 年夏天，Google DeepMind 公开了一段演示：桌上摆着几件玩具，其中有一只塑料恐龙；操作员给机械臂输入一句话——"把灭绝的动物捡起来"。机械臂伸向了恐龙。

这一幕值得多看几秒。训练这台机器人的示教数据里，没有一条示教叫"灭绝的动物"；"恐龙已经灭绝"这件事，是模型从互联网文本里学来的。这个叫 RT-2 的系统，把一个视觉--语言模型直接微调成了策略：图像和指令进去，动作出来，中间没有人手写一条规则。第11章"会看会说"的模型，第一次"会做"了。Hans Moravec 在 1980 年代就指出，让计算机下棋容易，让它像一岁小孩那样感知和运动却极难；RT-2 让不少人觉得，这道鸿沟也许能靠规模填平。
\storynote{Chelsea Finn}{斯坦福大学教授。2017 年与 Abbeel、Levine 提出元学习方法 MAML；她的实验室 2023 年提出的 ALOHA/ACT 让"动作块"成为 VLA 的标准配件。2024 年她与 Levine、Hausman 等人共同创立 Physical Intelligence，$\pi_0$ 出自该公司。}

于是一种说法流行起来：\textbf{有了大模型，机器人只要模仿人类示教就够了}——多采一些遥操作数据，多训几轮，剩下的交给规模。

这是本章要拆掉的第一个误解。模仿有三道天花板。第一道是\textbf{数据}：语言模型有整个互联网，机器人只有"一个人戴着设备做一遍才得到一条示教"，全世界的机器人数据加起来也不过是语料的零头。第二道是\textbf{示教者}：人会犹豫、会失误，模仿学到的至多是"平均水平的人"；要超过示教者，机器人得从自己的经验里学。第三道是\textbf{安全}：模仿的安全来自"数据里没有危险动作"，这是统计意义上的保证，不是数学意义上的——而第16章的 CBF 给出的恰恰是后者。三道天花板对应三个问题：数据从哪来、怎么超过示教者、怎么保证不出事。把它们翻译成本书的语言——目标是什么、约束是什么、$\lambda$ 代表什么——就是本章。

\vspace{0.3cm}
\noindent\textit{为什么一个会写诗的模型，让它"把杯子放进洗碗机"时还得重新学一门"语言"？为什么机器人一次要预测 50 步动作，而不是像下棋一样走一步看一步？为什么在仿真里撞一百万次墙不心疼，真机上撞一次就得停下来算账？}

\begin{quote}
\textbf{本章目标}：把"看图做事"的机器人写成一个约束优化问题，看清它的目标（模仿误差或回报）、约束（动力学、安全集、数据预算）与乘子（安全的价格、模型误差的价格、样本的价格）。你会看到，模仿训练的 VLA 是把约束交给数据"代入消元"、暂时不算 $\lambda$ 的做法；流匹配把"生成一段动作"变成"解一个 ODE"；世界--动作模型给"预测"与"行动"定相对价格；而真机强化学习，是全书的三个问题第一次必须同时回答的地方。读完本章，你应当能对任何一个具身智能系统追问：它在优化什么、它受什么约束、它的每一次犹豫是哪个 $\lambda$ 在起作用。
\end{quote}

\begin{tcolorbox}[colback=green!5, colframe=green!60!black, title=学完这章，你将能够...]
\begin{itemize}
    \item 把 VLA 的训练写成约束优化：目标（模仿或回报）、约束（动力学、安全集）与暂时缺席的乘子
    \item 比较三种动作表示——离散 token、连续动作块、压缩后离散——在精度、速度与训练难度上的权衡
    \item 解释流匹配为什么"训练是回归、推理是积分"，并用几步欧拉法从噪声生成一段动作块
    \item 说明世界--动作模型中的 $\lambda$ 为什么是"预测"与"行动"的相对价格，以及它与第14章 $\rho$ 的关系
    \item 列出真机强化学习的四个困难与五条路线，写出拉格朗日约束强化学习的对偶上升更新
    \item 分清"解出来的 $\lambda^*$"（CBF）与"学出来的 $\lambda$"（约束 RL）各给出什么样的安全保证
\end{itemize}
\end{tcolorbox}

本章分三步走：视觉--语言--动作模型（从"看图说话"到"看图做事"）$\to$ 世界--动作模型（预测与行动共用一个隐状态）$\to$ 真机强化学习（三条约束、三个 $\lambda$）$\to$ 结语：从牛顿到智能。

\begin{tcolorbox}[colback=yellow!10, colframe=orange!80, title=本章的核心观点]
\textbf{本章的问题仍是第1章的拉格朗日乘子法，只是变量换成了策略参数与动作块，而约束——动力学、安全集、数据预算——第一次同时到场。}模仿训练时它们被示教数据隐式满足，$\lambda$ 暂时缺席；一上真机做强化学习，三个 $\lambda$ 就得同时计算。

\begin{center}
\renewcommand{\arraystretch}{1.3}
\setlength{\tabcolsep}{3pt}
\footnotesize
\begin{tabular}{p{1.0cm}p{1.8cm}p{3.9cm}p{4.2cm}}
\toprule
\textbf{概念} & \textbf{第1章} & \textbf{本章：VLA（模仿）} & \textbf{本章：真机强化学习} \\
\midrule
变量 & $x$ & 策略参数 $\theta$，动作块 $a_{t:t+H}$ & 策略 $\pi$，与乘子 $\lambda$ 交替迭代 \\
目标 & $f(x)$ & 模仿误差 $-\log \pi_\theta(a \mid o, \ell)$，或流匹配的 $\|v_\theta - (a_1 - a_0)\|^2$ & 期望回报 $\E[\sum_t \gamma^t r_t]$ \\
约束 & $g = 0$，$h \le 0$ & 动力学 $x_{t+1} = f(x_t, a_t)$、安全集 $h(x) \ge 0$——都由示教数据隐式满足 & 动力学；安全集 $h \ge 0$ 或违约预算 $\E[\sum_t \gamma^t c_t] \le \delta$；真机步数 $\le N$ \\
乘子 & 影子价格 & 暂时缺席：约束被数据"代入消元" & 安全的价格 $\lambda$、模型误差的价格（伴随变量）、样本的价格 \\
一阶条件 & $\nabla f + \lambda \nabla g = 0$ & 回归：$v_\theta(a_\tau, \tau) \approx a_1 - a_0$ & 策略梯度上升 $+$ 对偶上升 $\lambda \leftarrow \max\{0,\ \lambda + \eta(\E[\sum_t \gamma^t c_t] - \delta)\}$ \\
影子价格 & 约束放松一单位，目标改善多少 & 多一条示教，模仿误差降多少 & 多容忍一次碰撞、多采一步真机数据，各换来多少回报 \\
\bottomrule
\end{tabular}
\end{center}
\end{tcolorbox}

% ====================================================================================
\section{视觉--语言--动作模型：从说话到做事}
% ====================================================================================

第16章谈的都是机器人的"内部"：世界模型在脑中模拟，CBF 在贴近边界时出手，可微物理把牛顿定律塞进反向传播。但一个机器人最终要面对的场景往往是这样的：桌上一堆杂物，有人对它说"把红色的杯子放进洗碗机"。它需要\textbf{看}（找到杯子和洗碗机）、\textbf{听}（理解这句话）、然后\textbf{做}（伸手、抓取、移动、放下）。看和听，第11章的视觉--语言模型已经会了；第12章又告诉我们，怎样用后训练——RLHF、DPO、可验证奖励——把一个预训练模型调教成"听话"的智能体。本章要把这两章的成果搬到机械臂上：先在 VLM 后面接一个动作头，让它"会做"；再问一句，后训练那一套到了真机上还灵不灵。
\review{第11章}{第11章的视觉--语言模型（VLM）把图像切成视觉 token，与文本 token 一起送进同一个 Transformer 联合训练——模型于是"会看会说"。本节要在它后面再接一个输出头，让它"会做"。}

第11章的视觉--语言模型（Vision-Language Model，VLM）已经解决了前两件事：它能看图，能回答"图里有什么、该先做什么"。缺的是第三件——\textbf{动作}。把动作也加进来，就得到\textbf{视觉--语言--动作模型}（Vision-Language-Action model，VLA）：
\begin{equation}
    \pi_\theta(a \mid o, \ell)
\end{equation}
其中 $o$ 是当前观测（一路或多路摄像头图像，加上关节角等本体感觉），$\ell$ 是语言指令，$a$ 是动作（例如机械臂末端的位姿增量，或各关节的目标角度）。

\subsection{策略也是一个约束优化}

用本书的语言，训练一个 VLA 就是在求解
\begin{equation}
    \min_\theta\ \underbrace{\E_{(o,\ell,a)\sim\mathcal{D}}\big[-\log \pi_\theta(a\mid o,\ell)\big]}_{\text{目标：模仿示教数据}} \quad \text{s.t.}\quad \underbrace{x_{t+1} = f(x_t, a_t)}_{\text{动力学}},\qquad \underbrace{h(x_t)\ge 0}_{\text{安全}}
    \label{eq:vla_problem}
\end{equation}
目标可以是\textbf{模仿}（示教数据的负对数似然，即第9章算法全景表里的 BC——行为克隆），也可以是\textbf{回报}（强化学习）；约束仍是第16章的那两条——物理动力学（等式）与安全集合（不等式）。此刻我们只用目标来训练：示教数据本身就是可行的（人做出来的动作当然满足动力学），安全约束则暂时靠"数据里没有危险动作"来隐式满足。这两条约束的乘子——它们各自的"价格"——留到下一节真机强化学习时再算。

\subsection{动作怎么表示：token、动作块与压缩}

VLM 的输出是文本 token，而机械臂需要的是连续的数值。中间怎么接？目前有三条路。

\paragraph{路线一：把动作离散成 token。}
最直接的做法是让动作"假装"成文字。RT-2（Google DeepMind，2023）把每个动作维度的数值量化成整数，再把整数当作词表里的 token，与文本共用同一个输出头——VLM 于是可以像续写句子一样"续写"动作。OpenVLA（2024）沿用了这个思路并开源：每个动作维度切成 256 个等宽的 bin，7 维动作就是 7 个 token。

代价有两条。\textbf{精度}：256 个 bin 对夹爪开合够用，对需要毫米级对准的插孔任务就粗了；量化误差是硬伤，模型再大也补不回来。\textbf{速度}：自回归解码一次只出一个 token，一个动作要跑 7 次前向（附录C：解码阶段受显存带宽限制）；要让机械臂以 50 Hz 控制，每 20 ms 就得出一个动作，几十亿参数模型的逐 token 解码根本跟不上。

\paragraph{路线二：连续动作加动作块。}
既然一次出一个动作太慢，就一次出一段：\textbf{动作块}（action chunking）让策略一次预测未来 $H$ 步的动作序列 $a_{t:t+H}$，即 $(a_t, a_{t+1}, \ldots, a_{t+H-1})$，执行完（或执行一部分后）再重新预测。这样每 $H$ 步只需推理一次；同时动作块是一个 $H \times d$ 的连续张量，可以用回归或生成模型（扩散、流匹配）直接输出，不必量化。这条路线的两块基石都出现在 2023 年：Tony Zhao 等人的 ALOHA 与 ACT（动作分块 Transformer）用两台廉价机械臂做双手精细操作，让"一次预测一段"成为标准做法；迟程（Cheng Chi）等人在哥伦比亚大学 Shuran Song 组提出的 Diffusion Policy，则用扩散模型来生成动作块——下面 $\pi_0$ 的流匹配沿的是同一条路。
\review{第7章}{MPC 是"规划 $N$ 步、只执行 1 步"；动作块是"规划 $H$ 步、执行 $k$ 步（$1 \le k \le H$）"。两者都在用"多看几步"换取平滑与远见，区别是 MPC 有显式模型，动作块的"模型"藏在策略网络里。}

动作块还有一个意外的好处：示教数据里人的动作常有停顿、犹豫和多模态（这次从左边绕，下次从右边绕）。逐步预测会把这些平均掉，得到一个"既不左也不右"的糟糕动作；预测整段则能把"一整条绕行路线"当作一个整体来学。生成模型正擅长表示这种多模态分布，这是扩散与流匹配被用来出动作块的根本原因。输入端同样有文章可做：许华哲实验室 2024 年提出的 3D Diffusion Policy（DP3，Yanjie Ze 等）用三维点云代替图像作为扩散策略的输入，把学会一个操作技能所需的示教数量压到了几十条——点云直接携带几何信息，策略不必再从像素里自己推断物体在哪里。

\paragraph{路线三：先压缩再离散。}
动作块是连续的，但 VLM 的自回归头只吃 token。$\pi_0$-FAST（2025）的折中是：先对动作块做\textbf{离散余弦变换}（DCT），把 $H \times d$ 的时间序列变成频域系数——平滑的动作序列高频分量接近零，可以扔掉——再把保留的系数量化、编码成少量 token。这和 JPEG 压缩图片是同一个思路：先变换、再量化、再熵编码。结果是同样一段动作块所需的 token 数大幅减少，自回归解码于是变得可以承受。

\subsection{$\pi_0$：用流匹配生成动作块}

Physical Intelligence 在 2024 年提出的 $\pi_0$ 是路线二的代表。它的结构分两块（图~\ref{fig:vla_arch}）。一块是 \textbf{VLM 主干}：一个预训练的视觉--语言模型（PaliGemma 一类），负责把图像、指令与本体状态编码成一串 token 表示——"看"和"听"在这里完成。另一块是\textbf{动作专家}（action expert）：一个较小的子网络，与主干共用注意力（动作 token 可以"看到"主干的 token），但有自己的一套权重；它不输出文本 token，而是输出一个连续的动作块（例如未来 50 步）。

训练分两阶段：先在\textbf{跨本体}（cross-embodiment）的大规模数据上预训练——数据来自多种不同的单臂、双臂平台，甚至移动操作机器人，动作维度不同就用零填充对齐；再在目标任务的高质量示教上做\textbf{后训练}（post-training）。这和第11章语言模型"预训练 + 微调"、第12章"后训练"的配方是一样的——只是这里后训练用的数据是示教，而不是人类偏好。

\begin{figure}[htbp]
\centering
\resizebox{0.98\textwidth}{!}{%
\begin{tikzpicture}[>=Stealth, node distance=0.4cm and 0.7cm,
    box/.style={draw, rounded corners, align=center, minimum height=1.0cm, font=\scriptsize},
    input/.style={box, fill=green!15, minimum width=2.0cm},
    vlm/.style={box, fill=blue!15, minimum width=2.4cm},
    expert/.style={box, fill=orange!20, minimum width=2.6cm},
    act/.style={box, fill=yellow!25, minimum width=1.6cm},
    robot/.style={box, fill=gray!20, minimum width=1.8cm}]
    \node[input] (img) {图像 $o_t$\\（多路相机）};
    \node[input, below=of img] (lang) {语言指令 $\ell$\\（"把杯子放进去"）};
    \node[input, below=of lang] (prop) {本体状态\\（关节角等）};
    \node[vlm, right=of lang] (vlm) {VLM 主干\\视觉 token 与文本 token\\联合注意力};
    \node[expert, right=1.3cm of vlm] (exp) {动作专家\\流匹配：从噪声 $a_0$ 出发\\对 $da/d\tau = v_\theta$ 积分到 $a_1$};
    \node[act, right=of exp] (chunk) {动作块\\ $a_{t:t+H}$};
    \node[robot, below=1.8cm of chunk] (rob) {机器人执行\\ $H$ 步（或前 $k$ 步）};
    \draw[->, thick] (img) -- (vlm);
    \draw[->, thick] (lang) -- (vlm);
    \draw[->, thick] (prop) -- (vlm);
    \draw[->, thick] (vlm) -- node[above, font=\tiny] {token 表示} (exp);
    \draw[->, thick] (exp) -- (chunk);
    \draw[->, thick] (chunk) -- (rob);
    \draw[->, thick, dashed] (rob.west) -| node[pos=0.25, below, font=\scriptsize] {新观测 $o_{t+H}$ 回到输入} ($(img.west)+(-0.6,0)$) -- (img.west);
    \begin{scope}[on background layer]
        \node[draw, dashed, rounded corners, fill=blue!3, inner sep=6pt, fit=(vlm)(exp), label={[font=\scriptsize]above:VLA 模型 $\pi_\theta(a_{t:t+H} \mid o_t, \ell)$}] {};
    \end{scope}
\end{tikzpicture}%
}
\caption{$\pi_0$ 一类 VLA 的结构。图像、语言指令与本体状态进入 VLM 主干，得到一串 token 表示；动作专家以此为条件，用流匹配从噪声积分出一个连续的动作块 $a_{t:t+H}$；机器人执行整块（或前 $k$ 步）后，新的观测回到输入，开始下一轮。虚线框内是被训练的策略 $\pi_\theta$。}
\label{fig:vla_arch}
\end{figure}

真正新的东西是动作专家怎么生成动作：它用的是\textbf{流匹配}（flow matching）。

\subsubsection{流匹配：生成一段动作 = 解一个 ODE}

我们想从条件分布 $p(a \mid o)$ 里采样（这里 $a$ 是整个动作块，$o$ 代表观测与指令）。流匹配的想法是：不直接学这个分布，而是学一个\textbf{速度场}，让噪声沿着它"流"成数据。

取 $a_0 \sim \mathcal{N}(0, I)$ 为噪声，$a_1$ 为数据集里的真实动作块。在两者之间画一条直线：
\begin{equation}
    a_\tau = \tau\, a_1 + (1 - \tau)\, a_0, \qquad \tau \in [0, 1]
\end{equation}
沿这条直线的速度是常数：$\frac{d a_\tau}{d\tau} = a_1 - a_0$。于是训练一个网络 $v_\theta(a_\tau, \tau \mid o)$ 去拟合它：
\review{第3章}{第3章用 E-L 方程证明"两点之间直线最短"。流匹配把噪声与数据之间的路径选成直线，取的正是这条"最省"的路径：它让待学的速度场尽量简单，推理时用很少的积分步就能走到头。}

\begin{keyformula}[title=流匹配：训练是回归，推理是积分]
\textbf{训练}（回归一个速度场）：
\begin{equation}
    \min_\theta\ \E_{\tau,\, a_0,\, (o, a_1)\sim\mathcal{D}} \Big\| v_\theta(a_\tau, \tau \mid o) - (a_1 - a_0) \Big\|^2, \qquad a_\tau = \tau a_1 + (1-\tau) a_0
\end{equation}
\textbf{推理}（解一个常微分方程）：从 $a_0 \sim \mathcal{N}(0, I)$ 出发，对
\begin{equation}
    \frac{d a}{d\tau} = v_\theta(a, \tau \mid o)
\end{equation}
从 $\tau = 0$ 积分到 $\tau = 1$（例如用 10 步欧拉法），终点就是生成的动作块。

\textbf{一句话}：训练时学"从噪声到动作该往哪走"，推理时把这个方向场积分一遍。生成一段动作 $=$ 解一个 ODE。
\end{keyformula}

为什么这个损失是对的？固定一对 $(a_0, a_1)$ 时，"正确速度"就是 $a_1 - a_0$；但同一个中间点 $a_\tau$ 可能来自很多不同的 $(a_0, a_1)$ 对，最小二乘回归学到的是它们的\textbf{条件期望}，而这个平均速度场恰好把边缘分布 $p(a_0)$ 输运到 $p(a_1)$。证明超出本书范围，但结论很干净：\textbf{逐样本的直线，平均起来是一个合法的速度场}。

\begin{insightbox}[title=流匹配与本书主线]
\textbf{推理是 ODE}。第7章把动力学 $\dot{x} = f(x, u)$ 当作约束，伴随变量是它的乘子；这里 $\frac{da}{d\tau} = v_\theta$ 是一条"人造的动力学"，$\tau$ 不是物理时间而是生成进度。生成一段动作，就是把一个 ODE 从噪声积分到数据。

\textbf{训练是回归}。目标函数是最小二乘，没有约束——"约束"全被直线插值吸收了。这和第8章"反向传播是逐层等式约束的乘子"形成对照：流匹配把约束消掉了，而不是给它配乘子；这正是第1章"代入消元"的做法。

\textbf{与扩散模型的关系}。扩散模型也是"噪声到数据"，但走的是弯曲的随机路径，推理往往要几十上百步；直线路径通常几步就够——这对需要高频控制的机器人至关重要。
\end{insightbox}

\subsection{代表性设计一览}

VLA 是一个变化很快的领域，表~\ref{tab:vla_compare}只记录几个有代表性的设计选择。重点看"动作输出"这一列——它决定了精度、速度与训练难度的权衡。GR00T N1（NVIDIA，2025）值得单独一提：它把"慢思考"与"快动作"明确拆成两个系统，VLM 负责低频的理解与规划，一个扩散 Transformer 负责高频地出动作——与 $\pi_0$ 的"主干 + 动作专家"殊途同归。
\think{自回归 token 也是"逐步"生成：每一步条件于之前所有 token。流匹配也是"逐步"：每一步更新整个动作块。这两种"逐步"沿的是哪两个不同的维度？为什么后者更适合高频控制？}

\begin{table}[htbp]
\centering
\small
\renewcommand{\arraystretch}{1.3}
\setlength{\tabcolsep}{4pt}
\begin{tabular}{p{2.0cm}p{0.9cm}p{3.4cm}p{5.3cm}}
\toprule
\textbf{模型} & \textbf{年份} & \textbf{动作输出} & \textbf{特点} \\
\midrule
RT-2 & 2023 & 离散 token，与文本共用词表 & 直接把 VLM 微调成策略；网页知识可迁移到操作任务 \\
OpenVLA & 2024 & 离散 token，每维 256 个 bin & 开源；在大规模跨本体数据集上训练 \\
$\pi_0$ & 2024 & 流匹配，连续动作块 & VLM 主干 + 动作专家；跨本体预训练 + 任务后训练 \\
$\pi_0$-FAST & 2025 & DCT 压缩后的离散 token & 自回归头也能输出动作块；token 数大幅减少 \\
$\pi_{0.5}$ & 2025 & 离散 token 预训练，流匹配后训练 & 网页、多本体、多环境的异构数据共训练；泛化到未见过的家庭环境 \\
GR00T N1 & 2025 & 扩散 Transformer，连续动作块 & 双系统：VLM 慢速理解 + 扩散头高频出动作 \\
\bottomrule
\end{tabular}
\caption{几个代表性 VLA 的设计选择。年份为论文或技术报告首次公开的年份。}
\label{tab:vla_compare}
\end{table}

\subsection{VLA 的三个未解之题}

\textbf{数据}。语言模型有整个互联网，VLA 只有遥操作：一个人戴着设备控制机器人做一遍，才得到一条示教。把全球几十个实验室的数据合在一起的跨本体数据集，也不过是语言模型语料的零头。跨本体预训练的动机正是"能凑的都凑上"；第12章的数据引擎（生成器--验证器--过滤器）在这里同样适用，只是"验证器"得在真实世界里给出判断——这正是下一节真机强化学习中"奖励"难题的来源；而下一节的世界--动作模型，则是想把没有动作标注的视频也用上。

\textbf{频率与延迟}。VLM 主干动辄数十亿参数，一次前向要几十到上百毫秒，而灵巧操作需要 50 Hz 上下的控制。动作块是目前的解药——推理一次、执行一段——但块内是\textbf{开环}的：这几十步里如果杯子被碰倒了，策略要等到下一次推理才知道。块越长，推理越省，反应越慢；这是一个反应速度与算力之间的权衡，没有免费午餐。

\textbf{安全没有保证}。VLA 的安全完全来自"训练数据里没有危险动作"，这是统计意义上的保证，不是数学意义上的：遇到一个没见过的场景，模型可能输出任何东西。对比第16章第~\ref{sec:cbf}~节：CBF 给出的是\textbf{前向不变}的证明——只要 $\dot{h} \ge -\alpha h$，就永远不会离开安全集合。一个自然的组合是把 VLA 的输出当作 $u_{\rm ref}$，外面套一层 CBF-QP 做安全滤波：学习负责"聪明"，约束负责"不出事"，$\lambda^*$ 负责在两者冲突时裁决。这正是下一节的起点。
\preview{下一节把两条线拧在一起：世界模型（预测）与 VLA（行动）合成一个模型；而"安全的价格"$\lambda$ 从 CBF 的闭式解，变成真机强化学习里随违约上升的对偶变量。}

% ====================================================================================
\section{展望：世界--动作模型与真机强化学习}
% ====================================================================================

第16章开头讲世界模型：给定动作，预测未来。上一节讲 VLA：给定观测，输出动作。两者是同一枚硬币的两面——一个是 $p(o_{t+1} \mid o_t, a_t)$，一个是 $\pi(a_t \mid o_t)$——却常常被当作两个模型分别训练。本节是"展望"：谈两个正在发生、尚未定型的方向，一个是把两面合成一枚硬币，一个是让机器人在真实世界里从自己的经验中学习。这一节的口吻会比前面谨慎，因为很多细节还在变；但它们背后的约束优化结构，现在就可以写清楚。

\subsection{世界--动作模型：预测与行动共用隐状态}

\paragraph{动机：视频是免费的监督。}
上一节末尾说过，VLA 的瓶颈是带动作标注的数据。但世界上有海量\textbf{没有}动作标注的视频——人做家务、拆快递、修自行车。这些视频里没有"关节角"，却有"世界怎么变"。如果模型除了输出动作，还被要求预测下一帧观测，那么预测下一帧本身就是一种监督：视频数据虽然不能教它"手该怎么动"，却能教它"手动了之后杯子会怎样"。

\textbf{世界--动作模型}（World Action Model，WAM）的思想就是把两个任务放进同一个模型里联合训练：
\begin{equation}
\begin{aligned}
    \min_\theta\ \ & \underbrace{\E\big[\|\hat{o}_{t+1} - o_{t+1}\|^2\big]}_{\text{预测未来观测}} + \lambda\, \underbrace{\E\big[\|\hat{a}_{t} - a_{t}\|^2\big]}_{\text{输出动作}} \\
    \text{s.t.}\ \ & z_t = E(o_t, \ell),\qquad \hat{o}_{t+1} = D(z_t, a_t),\qquad \hat{a}_t = \pi(z_t)
\end{aligned}
    \label{eq:wam}
\end{equation}
这里的约束说的是：预测头 $D$ 和动作头 $\pi$ 必须从\textbf{同一个}隐状态 $z_t$ 出发。没有这条约束，两项损失可以各学各的，等于两个独立模型；有了它，隐状态必须同时"够预测"和"够行动"——这正是第16章"潜空间 = 广义坐标"的要求，只是现在多了一位用户。而 $\lambda$ 是两项之间的权重：它决定了隐状态更偏向"看得准"还是"做得对"。（上式为了简洁写成了均方误差；实际系统里两项通常都是生成模型的损失——扩散、流匹配或 token 的交叉熵——但结构不变。）
\review{第14章}{PINN 损失 $L_{\rm PDE} + \rho L_{\rm BC}$ 里的 $\rho$ 也是两项之间的权重。第14章的教训：固定权重要靠试，自适应权重则按"哪项违约更多就给它涨价"的对偶上升来调。WAM 里的 $\lambda$ 面对的是同一个问题。}

\paragraph{几个方向。}
这类模型在 2025 年前后集中出现，路线不一。\textbf{统一世界模型}（UWM，2025）把视频扩散与动作扩散放进同一个 Transformer，两种模态各有自己的噪声进度：把动作那边的噪声开到最大，模型就退化成纯视频预测器，于是无动作标注的视频可以直接拿来训练；反过来把视频那边的噪声开到最大，它就是一个策略。\textbf{WorldVLA}（2025）用自回归 token 统一图像与动作：同一个模型既能"给定动作生成下一帧"，也能"给定图像生成动作"，两种任务互为对方的数据增强。\textbf{V-JEPA 2-AC}（2025）走的是第16章 LeCun 的路线：先在海量互联网视频上自监督地学一个视频表示（不预测像素，预测表征），再用少量机器人数据训练一个\textbf{动作条件}的预测器，规划时在表征空间里搜索动作序列——这是 Dreamer 式的"在梦中规划"，只是梦是在 JEPA 的表征空间里做的。这些设计谁会胜出，现在下结论为时过早。但它们都在回答同一个问题：如何让"预测"与"行动"共享一个表示，以及如何给这两项定价。

\begin{physicalmeaning}[title=WAM 里的三样东西]
\textbf{目标}：观测预测误差 $+$ 动作误差。\textbf{约束}：两者必须经过同一个隐状态。\textbf{$\lambda$}：两项的相对价格——它回答"为了让动作更准，我愿意让预测差多少"。这和第14章 PINN 里 $\rho$ 的角色完全一样，也一样需要自适应地调。
\end{physicalmeaning}

\subsection{真机强化学习为什么难}

无论是 VLA 还是 WAM，上面讲的都是\textbf{模仿}：从人的示教里学。模仿的天花板是示教者：人会犹豫、会失误，机器人学到的也就是"平均水平的人"。要超过示教者，就得让机器人从\textbf{自己的}经验中学——这就是第9章的强化学习。Rodney Brooks 1990 年在《大象不下棋》里主张的正是这一点：智能必须从与物理世界的交互中长出来，而不是先建一个符号模型。可是第9章的实验都在仿真里跑：仿真里一秒能走几千步，重置不要钱，撞墙也不心疼。到了真机上，四件事同时变难。

第一件是\textbf{样本贵}：每一步都是真实时间，仿真里一小时能采几百万步，真机上一小时也就几千步，而且机器人会磨损，电机会发热。视觉强化学习的样本效率因此成了一个独立的课题——许华哲实验室的 H-InDex（2023）用人手先验来做灵巧手的视觉强化学习，DrM（2024）则研究了视觉强化学习里神经元"休眠"导致的探索失败。第二件是\textbf{复位}：一次尝试结束后，谁把杯子放回原处？仿真里一行 \texttt{env.reset()}，真机上要么有人守着，要么得再训练一个"把东西放回去"的策略。

第三件是\textbf{安全}：探索意味着尝试没做过的动作，而没做过的动作里有一部分会撞坏东西——包括机器人自己；仿真里探索是免费的，真机上探索是有价格的。第四件是\textbf{奖励}：仿真器知道杯子的精确位置，能算出"离目标还有多远"；真机上只有摄像头图像，"成功了吗"本身就得靠一个分类器来判断，而分类器也会被策略"钻空子"。

\subsection{五条路线}

这四个困难没有一个通用解法，目前的实践是几条路线叠加。
\storynote{Sergey Levine}{加州大学伯克利分校教授，把深度强化学习用于机器人的主要推动者之一，SERL、HIL-SERL 出自他的实验室；2024 年他与同事共同创立了 Physical Intelligence，$\pi_0$ 系列即出自该公司。}

\paragraph{(a) 先在仿真里学，再搬到真机。}
Sim-to-real 的做法是在仿真里做全部探索，然后把策略迁移到真机；为了让策略不依赖仿真的具体参数，训练时随机扰动质量、摩擦、光照、相机位置——\textbf{域随机化}（domain randomization）——逼策略对"仿真与真实之间的差距"鲁棒。第16章"可微物理"一节的系统辨识是它的另一半：先用真机数据把仿真器校准得更像真实世界，再在里面训练。四足机器人行走是这条路线最成功的应用之一；灵巧操作则因为接触动力学难以仿真，迁移仍然困难。

\paragraph{(b) 离线起步，在线微调。}
第9章讲过离线强化学习：不与环境交互，只从固定数据集学。真机 RL 的实用配方是先用示教数据做行为克隆或离线 RL，得到一个"能用"的初始策略，再在真机上做少量在线交互微调——探索从一个不错的起点出发，样本贵的问题就轻了很多。

\paragraph{(c) 人在回路的真机强化学习。}
SERL（2024）和 HIL-SERL（2024）把上面的想法做成了一套可复用的系统：示教起步、离策略 RL（经验池里同时放示教与自身经验，第9章的 Off-policy）、\textbf{人在回路}——训练过程中人可以随时接管遥操作，纠正一段动作，纠正的片段也进经验池；再加一个从图像判断成败的奖励分类器。结果是几小时的真机训练就能把精密插装、翻转物体、装配这类任务的成功率提到接近百分之百，而且策略的动作比示教者更快。人类纠正在这里的角色很有意思：它既是探索的引导（省样本），也是安全网（防撞），还是奖励的补充（"我接管了"本身就是一种负反馈）。

\paragraph{(d) 对 VLA 做强化学习后训练。}
上一节的 VLA 是用模仿训练的。既然它已经"能用"，就可以把它当作路线 (b) 的起点，用真机经验继续改进。Physical Intelligence 在 2025 年的 $\pi^{*}_{0.6}$ 用一种叫 RECAP 的方法走了这条路，思想大致是：让机器人自己执行任务，收集成功与失败的经验，加上人类的纠正，训练一个价值函数估计"这个状态离成功还有多远"；策略不是直接最大化回报，而是以\textbf{优势}（advantage，第9章 PPO 里的 $A = Q - V$）为条件来训练——训练时告诉它"这段动作比平均好多少"，推理时要求它输出"比平均好"的动作。好处是不必改动 VLA 的架构与损失，只是输入里多了一个条件；思想上它是第9章策略改进的一种"温和"版本。细节仍在演进，这里只概述思路。

\paragraph{(e) 把安全写进约束。}
以上四条路线主要在解决样本、复位与奖励；安全需要单独处理。第16章 CBF 一节给出了一种答案：安全是不等式约束 $h(x) \ge 0$，进入 KKT 后其乘子 $\lambda^*$ 是"安全的价格"，只在贴近边界时非零；把它套在 RL 策略外面做安全滤波，探索就被限制在安全集合内。另一种答案是\textbf{拉格朗日约束强化学习}：把"期望违约代价不超过预算 $\delta$"写成约束，
\begin{equation}
    \max_\pi\ \E\Big[\sum_t \gamma^t r_t\Big] \quad \text{s.t.}\quad \E\Big[\sum_t \gamma^t c_t\Big] \le \delta,
\end{equation}
其中 $c_t \ge 0$ 是第 $t$ 步的违约代价（例如碰撞记 1）。写出拉格朗日函数 $\Lag(\pi, \lambda) = \E[\sum_t \gamma^t r_t] - \lambda\big(\E[\sum_t \gamma^t c_t] - \delta\big)$，对 $\pi$ 做策略梯度上升，对 $\lambda$ 做对偶上升：
\begin{equation}
    \lambda \leftarrow \max\Big\{0,\ \lambda + \eta\Big(\E\Big[\sum_t \gamma^t c_t\Big] - \delta\Big)\Big\}.
\end{equation}
违约越多，$\lambda$ 涨得越快，惩罚越重；不违约时 $\lambda$ 降回零，策略专心追求回报。
\review{第2章}{对偶上升 $\lambda \leftarrow \lambda + \eta g$：乘子沿"约束违反量"的方向走，是第2章对偶问题的梯度法，第14章的自适应权重用的是同一条更新式。这里的 $g$ 是违约代价超出预算的部分，$\max\{0,\cdot\}$ 保证不等式约束的乘子非负。}

注意两种答案的分工：CBF 的 $\lambda^*$ 是每一时刻在 QP 里\textbf{解出来}的，给出逐步的硬保证；约束 RL 的 $\lambda$ 是在训练过程中\textbf{学出来}的，给出的是期望意义下的软保证。前者的代价是需要一个动力学模型和一个障碍函数 $h$，后者的代价是训练期间仍会违约——这正是为什么真机上两者常常叠用。

\subsection{收束：三条约束，三个 $\lambda$}

回头看，真机强化学习恰好是本书框架的终点：它把全书出现过的约束一次性全部请到场上。

\begin{keyidea}[真机强化学习：本书框架的终点]
\textbf{目标}：期望回报（或成功率），必要时加上模仿项以稳定起步。

\textbf{约束}：三条同时在场。\textbf{物理动力学} $x_{t+1} = f(x_t, a_t)$ 是等式约束（第5、7章），乘子是伴随变量 / 值函数梯度，回答"状态多动一点，回报变多少"；\textbf{安全集合} $h(x_t) \ge 0$ 是不等式约束（第16章 CBF 一节），乘子是安全的价格，只在贴近边界时非零；\textbf{数据与时间预算}——真机交互步数 $\le N$，人类干预次数 $\le M$——的乘子是"再多采一步真机数据值多少回报"，仿真预训练、离线起步、动作块、人类纠正，全都是在压低这个价格。

\textbf{$\lambda$}：三个乘子同时在算。它们的量纲都是"每放松一单位约束换来多少回报"，所以可以互相比较——当安全的价格远高于数据的价格时，就该多花数据换安全；反之亦然。全书的三个问题，在这里第一次要同时回答。
\end{keyidea}

一个机器人在真实世界里学习，就是在这三条约束之下，用有限的经验寻找更好的策略；而它每一次犹豫、每一次被人接管、每一次贴着边界绕行，都是某个 $\lambda$ 在起作用。

到这里，全书的三个问题——目标是什么、约束是什么、$\lambda$ 代表什么——第一次在同一个系统里被同时问出，也被同时回答。小结与习题之后，本书的结语会把这条从牛顿到智能的线从头再走一遍：你会发现，两百多年前为力学发明的那个乘子，从未离开过舞台。


% ====================================================================================
\section{本章小结}
% ====================================================================================

\begin{enumerate}
    \item \textbf{VLA}：$\pi_\theta(a \mid o, \ell)$ 在第11章 VLM 的"会看会说"后面接一个动作头，让它"会做"。训练写成约束优化 \eqref{eq:vla_problem}：目标是模仿误差（或回报），约束是动力学与安全集；模仿训练时两条约束都由示教数据隐式满足，$\lambda$ 暂时缺席。
    \item \textbf{动作表示三条路}：离散成 token（RT-2、OpenVLA：量化误差与逐 token 解码的延迟是硬伤）；连续动作块（一次预测 $H$ 步、执行 $k$ 步，是 MPC"规划多步、只走一步"的近亲，还顺手避免了示教的多模态被平均掉）；先 DCT 压缩再离散（$\pi_0$-FAST）。
    \item \textbf{流匹配}：训练是回归——在噪声 $a_0$ 与数据 $a_1$ 之间画直线，学速度场 $v_\theta \approx a_1 - a_0$；推理是积分——对 $da/d\tau = v_\theta$ 从 $\tau = 0$ 走到 $\tau = 1$。直线是第3章"最省"的路径；约束被插值吸收，是第1章的"代入消元"。
    \item \textbf{VLA 的三个未解之题}：数据只有遥操作；VLM 主干慢、动作块内开环，反应速度与算力之间没有免费午餐；安全只有统计意义的保证——把输出当作 $u_{\rm ref}$ 套一层 CBF-QP，是"学习负责聪明、约束负责不出事"的自然分工。
    \item \textbf{世界--动作模型}：预测头与动作头必须从同一个隐状态出发，这条约束让没有动作标注的视频也成为监督；$\lambda$ 是两项损失的相对价格，与第14章 PINN 的 $\rho$ 面对同一个"权重怎么定"的问题。
    \item \textbf{真机强化学习的四难五路}：样本、复位、安全、奖励四件事同时变难；路线是仿真迁移（域随机化加第16章的系统辨识）、离线起步在线微调、人在回路（HIL-SERL）、对 VLA 做 RL 后训练（以优势为条件）、把安全写进约束。
    \item \textbf{两种安全的 $\lambda$}：CBF 的 $\lambda^*$ 每一时刻在 QP 里解出来，给逐步的硬保证；拉格朗日约束 RL 的 $\lambda$ 在训练中由对偶上升学出来，给期望意义下的软保证。真机上两者常常叠用。
    \item \textbf{三条约束，三个 $\lambda$}：动力学、安全集、数据预算同时在场，三个乘子的量纲都是"每放松一单位约束换来多少回报"，因而可以互相比较——全书三问在这里第一次要同时回答。
\end{enumerate}

\begin{physicalmeaning}
\textbf{$\lambda$ 在本章的意义}：

\textbf{安全约束的价格}。同一条不等式约束 $h \ge 0$，有两种算法算它的乘子。CBF 把它写进 QP，$\lambda^*$ 由互补松弛闭式定出——远离边界为零，贴近边界才非零（第1章）；约束 RL 把它当作对偶变量，按 $\lambda \leftarrow \max\{0,\ \lambda + \eta(\E[\sum_t \gamma^t c_t] - \delta)\}$ 随违约上升、不违约时回落——这是沿第2章对偶函数的梯度做上升，也是第14章自适应权重的同一条更新式。

\textbf{模型误差的价格}。动力学约束 $x_{t+1} = f(x_t, a_t)$ 的乘子是第7章的伴随变量、第9章的值函数梯度，它回答"模型在这一步错一点，回报差多少"。价格高的地方该多花真机数据做系统辨识（第16章），价格低的地方尽管信仿真；WAM 里的 $\lambda$ 是它的训练版——为了动作更准，愿意让预测差多少。

\textbf{样本预算的价格}。"真机步数 $\le N$"的乘子是"再多采一步真机数据值多少回报"。仿真预训练、离线起步、动作块、人类纠正，全都是在压低它。第12章 RLHF 里 KL 约束的乘子 $\beta$ 说的是同一件事：策略离参考模型（这里是示教策略）能走多远，是有价格的——走得越远，样本越贵，也越危险。

回头看：第16章里 CBF 的 $\lambda^*$ 是安全的斥力，可微物理的 $\lambda_t$ 是伴随变量，自由能的乘子是惊奇；本章把它们请到同一台机器人身上，让三个价格互相比较。从第1章的影子价格到这里，全书自始至终只有一个 $\lambda$。
\end{physicalmeaning}

% ====================================================================================
\section{习题}
% ====================================================================================

\paragraph{计算题}
\begin{enumerate}
    \item \textbf{动作 token 的精度与速度}：
    某 VLA 把每个关节角离散成 256 个等宽的 bin，关节角范围为 $[-\pi, \pi]$；机械臂从基座到末端长 $0.8$ m，动作有 7 维，一次前向推理耗时 30 ms。
    \begin{enumerate}
        \item 一个 bin 对应多大的角度？若该关节位于基座，末端位置的量化误差上界大约是多少？这对"把插头插进插座"（公差约 1 mm）意味着什么？
        \item 若动作逐 token 自回归解码，出一个动作要多久？能支持多高的控制频率？与 50 Hz 的需求相差几倍？
        \item 改为一次推理输出 $H = 50$ 步的连续动作块，推理耗时升到 100 ms。平均到每一步的推理时间是多少？若只执行前 $k = 25$ 步就重新推理呢？
    \end{enumerate}

    \item \textbf{一维流匹配的"平均陷阱"}：
    设数据 $a_1$ 以等概率取 $-1$ 或 $+1$（示教里的"从左绕"与"从右绕"），噪声 $a_0 \sim \mathcal{N}(0, 1)$，直线插值 $a_\tau = \tau a_1 + (1 - \tau) a_0$。
    \begin{enumerate}
        \item 在 $\tau = 0$ 处，最优速度场是 $v^*(a, 0) = \E[a_1 - a_0 \mid a_0 = a]$。算出它。
        \item 用一步欧拉法（步长 1）从 $a_0 = a$ 出发，落在哪里？这对应正文里"逐步预测把多模态平均掉"的哪种失败？
        \item 当 $\tau$ 接近 1 时，$a_\tau$ 靠近 $+1$ 的点，其 $a_1$ 更可能是哪一个？据此说明为什么多步积分能把两个模态分开。
    \end{enumerate}

    \item \textbf{对偶上升的数值演算}：
    约束 RL 中违约预算 $\delta = 0.1$，步长 $\eta = 0.5$，$\lambda_0 = 0$。前三轮策略的期望违约代价依次为 $0.3,\ 0.3,\ 0.2$；之后策略在 $\lambda$ 的惩罚下把违约压到 $0.05$ 并保持。
    \begin{enumerate}
        \item 逐轮写出 $\lambda_1, \lambda_2, \lambda_3$，以及违约降到 $0.05$ 之后再过四轮的 $\lambda$。
        \item 若策略无论如何都做不到违约 $\le \delta$（问题不可行），$\lambda$ 会怎样？从对偶的角度解释这个现象，并说明实践中如何发现它。
    \end{enumerate}
\end{enumerate}

\paragraph{证明题}
\begin{enumerate}
    \setcounter{enumi}{3}
    \item \textbf{约束 RL 的弱对偶}：
    记 $J_r(\pi) = \E[\sum_t \gamma^t r_t]$，$J_c(\pi) = \E[\sum_t \gamma^t c_t]$，$\Lag(\pi, \lambda) = J_r(\pi) - \lambda\,(J_c(\pi) - \delta)$，$d(\lambda) = \max_\pi \Lag(\pi, \lambda)$。
    \begin{enumerate}
        \item 证明：对任意 $\lambda \ge 0$ 与任意满足 $J_c(\pi) \le \delta$ 的策略 $\pi$，都有 $d(\lambda) \ge J_r(\pi)$；因此 $\min_{\lambda \ge 0} d(\lambda)$ 是原问题最优值的上界（第2章的弱对偶）。
        \item 设 $\pi_\lambda$ 是 $\Lag(\cdot, \lambda)$ 的最大化者。证明 $d$ 在 $\lambda$ 处的（次）梯度是 $-(J_c(\pi_\lambda) - \delta)$，从而对 $d$ 做梯度下降恰是正文里的对偶上升式。
        \item 若最优处 $\lambda^* > 0$，由互补松弛说明违约预算一定被用满：$J_c(\pi^*) = \delta$。用"价格"的语言解释这意味着什么。
    \end{enumerate}

    \item \textbf{流匹配的最小二乘解}：
    正文说流匹配学到的是"平均速度场"。证明：对固定的 $(a, \tau)$，使 $\E\big[\|v - (a_1 - a_0)\|^2 \,\big|\, a_\tau = a\big]$ 最小的 $v$ 是条件期望 $v^*(a, \tau) = \E[a_1 - a_0 \mid a_\tau = a]$。再说明：若每个 $a_\tau$ 只对应唯一一对 $(a_0, a_1)$（例如 $a_1$ 是 $a_0$ 的确定函数），则沿 $v^*$ 积分恰好从 $a_0$ 走到 $a_1$，一步欧拉法就够了。
\end{enumerate}

\paragraph{编程题}
\begin{enumerate}
    \setcounter{enumi}{5}
    \item \textbf{玩具流匹配}：
    实现计算题 2 的一维例子：数据 $a_1 \in \{-1, +1\}$（可加少量高斯噪声），用一个小 MLP $v_\theta(a, \tau)$ 拟合 $a_1 - a_0$；训练后分别用 1、2、5、20 步欧拉法从 $\mathcal{N}(0, 1)$ 采样，画出样本直方图。观察一步积分是否坍缩到均值附近、多少步才能分出两个模态。再把 $a_1$ 换成二维平面上的两条"绕行路线"（各 10 个路点），验证动作块能整体地保留多模态。

    \item \textbf{学习策略加 CBF 滤波}：
    在 \texttt{code\_chap16/chap16\_world\_model.py} 的 CBF 避障实验上，把标称控制 $u_{\rm ref}$ 换成一个"探索中的策略"：在朝向目标的方向上叠加零均值高斯扰动，标准差 $\sigma$ 模拟探索的大胆程度。对 $\sigma = 0, 0.2, 0.5, 1.0$ 分别统计 $\lambda^* > 0$ 的时间比例与到达目标的平均步数，画出两者随 $\sigma$ 的曲线。解释：探索越大胆，安全约束的"价格"为什么越高？这与正文"样本的价格与安全的价格可以互相比较"有什么关系？
\end{enumerate}

\paragraph{思考题}
\begin{enumerate}
    \setcounter{enumi}{7}
    \item \textbf{VLA 与安全滤波}：
    设一个 VLA 策略每次输出 $H = 50$ 步的动作块 $a_{t:t+50}$，机器人以 50 Hz 执行，一次推理耗时 $T_{\rm inf}$。
    \begin{enumerate}
        \item 块内是开环的：如果执行到第 10 步时杯子被碰倒，策略最早在什么时候能"看到"这件事？用 $H$ 与 $T_{\rm inf}$ 写出反应延迟的上界，并说明只执行前 $k < H$ 步再重新推理会如何改变这个上界。
        \item 现在把每一步动作先送进 CBF-QP 做安全滤波（把 VLA 的输出当作 $u_{\rm ref}$）。$\lambda^*$ 在什么情况下非零？这时机器人执行的还是 VLA "想"做的动作吗？
        \item 流匹配的推理从噪声 $a_0$ 沿 $da/d\tau = v_\theta$ 积分到 $a_1$。如果在积分的每一步都做一次 CBF 投影，得到的还是原分布的样本吗？这与"先生成、再滤波"有什么区别？（提示：想想约束是加在"生成进度 $\tau$"上，还是加在"物理时间 $t$"上。）
    \end{enumerate}

    \item \textbf{人类纠正是哪个 $\lambda$}：
    HIL-SERL 里，人可以随时接管遥操作。正文说这一个动作同时扮演了三个角色：探索的引导、安全网、奖励的补充。
    \begin{enumerate}
        \item 把"人类干预次数 $\le M$"写成一条约束，它的乘子是什么量纲？什么时候值得多花一次人力？
        \item 用三条约束（动力学、安全集、数据预算）的乘子，分别解释人类纠正的三个角色。
        \item 若把人类纠正的片段以更高的权重放进经验池，这相当于在给哪条约束"涨价"？
    \end{enumerate}

    \item \textbf{RLHF 的 $\beta$ 与真机 RL 的 $\lambda$}：
    第12章说过，RLHF 目标 $\E[r] - \beta\,\mathrm{KL}(\pi \,\|\, \pi_{\rm ref})$ 里的 $\beta$ 就是 KL 约束的乘子。对 VLA 做 RL 后训练时，$\pi_{\rm ref}$ 是模仿训练出来的策略。
    \begin{enumerate}
        \item 把它改写成"$\mathrm{KL} \le \epsilon$"的约束形式，写出更新 $\beta$ 的对偶上升式，并与正文的约束 RL 对照：KL 约束与违约预算约束分别在给什么定价？
        \item 真机上"离示教策略走得太远"会同时触发哪两个价格上涨？这能解释为什么真机 RL 几乎都从示教起步吗？
        \item RECAP 不直接最大化回报，而是以优势为条件训练。从"约束"的角度说说这种做法为什么更温和：它让哪一条约束几乎自动满足了？
    \end{enumerate}
\end{enumerate}

% ====================================================================================
\section*{结语：从牛顿到智能}
\addcontentsline{toc}{section}{结语：从牛顿到智能}
% ====================================================================================
\vspace{0.5cm}

当牛顿在1687年写下$F = ma$时，他大概不会想到，这个简单的方程会成为理解智能的钥匙。

\vspace{0.3cm}

当拉格朗日在1788年用变分法重写力学时，他大概不会想到，他的$\lambda$乘子会在两百多年后被用来训练神经网络、控制机器人、理解大脑。

\vspace{0.3cm}

而今天，当我们站在人工智能革命的浪潮中，回望这本书走过的旅程——从\textbf{第1章}的KKT条件到\textbf{第9章}的强化学习，从\textbf{有限元}的$KU=F$到\textbf{PINN}的物理信息损失，从\textbf{单智能体}的最优控制到\textbf{多智能体}的纳什均衡，从\textbf{经典力学}的约束力到\textbf{世界模型}的潜空间动力学——我们发现，所有这些看似迥异的领域，竟然都在诉说同一个故事：
\begin{keyformula}[title=全书的核心信息]
\begin{center}
\Large
\textbf{智能 = 在约束下寻找最优}
\end{center}
\vspace{0.3cm}
无论是石头滚下山坡，还是AlphaGo落子博弈；无论是单摆的简谐振动，还是GPT生成流畅的文字——它们都在做同一件事：在各自的约束条件下，寻找某种意义上的"最优解"。
\vspace{0.3cm}
\textbf{约束}定义了可能性的边界，\textbf{优化}选择了边界内的最佳路径。
\textbf{$\lambda$}不只是数学技巧，它是约束的"价格"，是物理的"力"，是信息的"代价"。
\end{keyformula}
\vspace{0.5cm}
\begin{philosophy}
\textbf{给读者的临别赠言}
\vspace{0.3cm}
如果你是大二的学生，正站在科学与工程的分岔路口，希望这本书能给你一个统一的视角：不要把物理、数学、计算机看成孤立的学科，它们是同一座大厦的不同房间。
\vspace{0.3cm}
当你下次遇到一个复杂问题时，不妨问自己三个问题：\textbf{目标是什么？}——我想最大化或最小化什么；\textbf{约束是什么？}——什么条件必须满足；\textbf{$\lambda$代表什么？}——约束的代价、价格或力是什么。这三个问题，能帮你把任何复杂问题，翻译成你已经熟悉的数学语言。
\end{philosophy}